\begin{enumerate}
		\item \begin{enumerate}
			\item En 2026, la production diminue de 5~\% par rapport à celle de 2025, elle est donc multipliée par \np{0.95}.

			$\np{3000} \times 0,95 = 3 \times 1000 \times 0,95 = 3 \times 950 = 3\times (\np{1000} - 50) = \np{3000} - 150 = \np{2850}$.

			La production en 2026 sera donc de \np[~kWh]{2850}.

			\item Diminuer une quantité de 5~\% revient à multiplier celle-ci par \np{0.95}.

			Puisque la production de l'année $n+1$ baisse de 5~\% par rapport à celle de l'année $n$, on a donc, pour tout entier naturel  \quad \mbox{$n, u_{n+1} = \np{0.95} \times u_n$}.

			\item La relation de récurrence de la suite $(u_n)$ est celle d'une suite géométrique de raison \np{0.95}.
		\end{enumerate}

		\item \begin{enumerate}
			\item On lit dans le tableau (ligne 24) que la valeur de $u_n$ est inférieure à \np{1000} lorsque $n$ devient supérieur ou égal à $22$.

			$2025 + 22 = 2047$.

			C'est donc à partir de 2047 que la production sera inférieure à \np[~kWh]{1000}.

			\item D'après le tableau (ligne 27), la production sera de \np[~kWh]{832.17} au bout de 25 années, soit environ 28~\% de la quantité initiale, donc l'affirmation du fabricant est fausse.
		\end{enumerate}

		\item Le graphique 1 ne peut pas représenter $(u_n)$. En effet, les points sont alignés, ce qui caractérise une suite arithmétique, alors que $(u_n)$ est une suite géométrique.

		C'est donc le graphique 2 qui représente $(u_n)$.
	\end{enumerate}
